% Part: sets-functions-relations
% Chapter: arithmetization
% Section: From N to Z

\documentclass[../../../include/open-logic-section]{subfiles}


\begin{document}

\olfileid[fr]{sfr}{arith}{int}
\olsection{De $\Nat$ à $\Int$}

Partons de deux observations simples :
	\begin{enumerate}
		\item Tout entier relatif s'écrit sous la forme $n - m$, avec $n, m \in\Nat$.
		\item L'information contenue dans l'expression $n - m$ peut tout aussi bien être représentée par le couple $\tuple{n, m}$.
	\end{enumerate}
Nous savons déjà que les couples de nombres naturels sont les !!{element}s de $\Nat^2$, et nous supposons connue la structure de $\Nat$. Ces deux observations suggèrent donc une première idée, un peu naïve : \emph{prendre les entiers relatifs pour des couples de nombres naturels}.

Cette idée est en fait trop naïve. Nous voulons bien entendu que $0- 2 = 4 - 6$. Or $\tuple{0, 2 }\neq \tuple{4, 6}$. Nous ne pouvons donc pas simplement déclarer que $\Nat^2$ est l'ensemble des entiers relatifs.

En généralisant le problème précédent, voici la propriété recherchée :
	$$a - b = c - d \liff a + d = c + b$$
(C'est bien ainsi que les entiers relatifs sont \emph{censés} se comporter : il suffit d'ajouter $b$ et $d$ aux deux membres.) Pour garantir cette propriété, définissons sur les couples une relation d'équivalence $\Intequiv$ par :
$$\tuple{ a, b } \Intequiv \tuple{c, d}\liff a + d = c + b$$
Il faut maintenant vérifier qu'il s'agit bien d'une relation d'équivalence.
\begin{prop} $\Intequiv$ est une relation d'équivalence.
\end{prop}
	\begin{proof}
	Nous devons montrer que $\Intequiv$ est réflexive, symétrique et transitive.
	
	\emph{Réflexivité :} Nous avons $\tuple{a, b} \Intequiv \tuple{a, b}$ ; en effet, $a + b = b + a$.
	
	\emph{Symétrie :} Supposons que $\tuple{a, b} \Intequiv \tuple{c, d}$, c'est-à-dire que $a + d = c + b$. Alors $c + b = a + d$, donc $\tuple{c, d} \Intequiv \tuple{a, b}$.
	
	\emph{Transitivité :} Supposons que $\tuple{a, b} \Intequiv \tuple{c, d}\Intequiv \tuple{m, n}$. Nous avons donc $a + d = c + b$ et $c + n = m + d$. Par addition, $a + d + c + n = c + b + m + d$, d'où $a + n = m + b$ par simplification dans $\Nat$. Ainsi, $\tuple{a, b } \Intequiv \tuple{m, n}$.
\end{proof}

Nous pouvons à présent former les classes d'équivalence pour cette relation :

\begin{defn}
		Les entiers relatifs sont les classes d'équivalence, pour $\Intequiv$, de couples de nombres naturels ; autrement dit, $\Int = \equivclass{\Nat^2}{\Intequiv}$.
\end{defn}

Cette définition, posée par convention, peut susciter des réactions \emph{philosophiques} très diverses. Avant de les examiner, poursuivons toutefois l'étude de quelques points techniques.

Après avoir défini les entiers relatifs, nous devons définir les fonctions et les relations usuelles sur cet ensemble. Notons $\equivrep{m, n}{\Intequiv}$ la classe d'équivalence pour $\Intequiv$ qui contient le couple $\tuple{m, n}$ comme !!{element}.\footnote{Avec la notation introduite dans la \olref[sfr][rel][eqv]{def:equivalenceclass}, nous écririons $\equivrep{\tuple{m,n}}{\Intequiv}$ pour désigner le même objet. La notation abrégée est simplement plus lisible.} Ainsi :
	$$\equivrep{m, n}{\Intequiv} = \Setabs{\tuple{a, b} \in \Nat^2}{\tuple{a, b}\Intequiv \tuple{m, n}}$$
Posons maintenant :
	\begin{align*}
		\equivrep{a, b}{\Intequiv} + \equivrep{c, d}{\Intequiv} &= \equivrep{a + c, b + d}{\Intequiv}\\
		\equivrep{a, b}{\Intequiv} \times \equivrep{c, d}{\Intequiv} &= \equivrep{a c + b  d, a  d + b c}{\Intequiv}\\
		\equivrep{a, b}{\Intequiv} \leq \equivrep{c, d}{\Intequiv} &\liff a + d \leq b + c
	\end{align*}
(Comme il est d'usage, nous écrivons $ab$ pour $(a \times b)$ afin d'alléger les axiomes.) Il reste à nous assurer que ces définitions ont bien les propriétés \emph{attendues}. Préciser ces propriétés et les vérifier demande un travail assez long ; les détails sont reportés à la \olref[check]{sec}. Retenons pour l'instant que tout fonctionne !{}

Il reste un dernier point. Nous avons construit les entiers relatifs à partir
des nombres naturels. Dans cette construction, les nombres naturels
\emph{ne sont pas eux-mêmes des entiers relatifs}.\footnote{Précision éditoriale : il s'agit de distinguer les objets de départ des classes que l'on vient de construire, avant l'identification naturelle décrite ci-dessous ; cette phrase ne pose pas une impossibilité d'inclusion pour toute représentation ensembliste des nombres.}
Nous reviendrons sur la portée philosophique de cette distinction dans la
\olref[ref]{sec}. Du point de vue technique, il nous faut cependant un moyen
de traiter les nombres naturels \emph{comme} des entiers relatifs. L'idée
est simple : pour chaque $n \in \Nat$, posons
$n_\Int = \equivrep{n, 0}{\Intequiv}$. Il faut vérifier que cette définition
respecte les opérations et l'ordre, c'est-à-dire que, pour tous $m, n \in \Nat$,
\begin{align*}
	(m + n)_\Int &= m_\Int + n_\Int\\
	(m \times n)_\Int &= m_\Int \times n_\Int\\
	m \leq n &\liff m_\Int \leq n_\Int
\end{align*}
Ces vérifications sont immédiates. Pour établir la deuxième identité,
par exemple, il suffit d'utiliser les propriétés des nombres naturels :
%\begin{proof}
%	En utilisant les propriétés des nombres naturels, nous avons :
	\begin{align*}
%		(m + n)_\Int  = \equivrep{m + n, 0}{\Intequiv} = \equivrep{m + n, 0 + 0}{\Intequiv} = \equivrep{m, 0}{\Intequiv} + \equivrep{n, 0}{\Intequiv} = m_\Int + n_\Int\\
		(m \times n)_\Int  &= \equivrep{m \times n, 0}{\Intequiv} \\
		&= \equivrep{m \times n + 0 \times 0, m \times 0 + 0 \times n}{\Intequiv} \\
		&= \equivrep{m, 0}{\Intequiv} \times \equivrep{n, 0}{\Intequiv} \\
		&= m_\Int \times n_\Int
	\end{align*}
%\end{proof}
La vérification des deux autres conditions est laissée en exercice.
\begin{prob}
	Montrer que $(m + n)_\Int = m_\Int + n_\Int$ et que $m \leq n \liff m_\Int \leq n_\Int$, pour tous $m, n \in \Nat$.
\end{prob}
\end{document}
